%! AP Statistics — Unit 8, Lesson 8.4 — Homework (Answer Key)
\documentclass[10pt]{article}
\usepackage{apstats-article}
\usepackage{apstats-key}   % pulls in -boxes; do NOT also load -boxes
\usepackage{grading}
\gradingassignment{apstats-u08-l04-hw}
\gradingcourse{advanced_statistics}
\newcommand{\TallMath}[1]{$\displaystyle #1\rule[-1.4em]{0pt}{3.2em}$}

\begin{document}

\pageheader{Unit 8, Lesson 8.4}{Homework --- Answer Key}

\vspace{0.1in}
\textbf{Directions:} Show all work. Express expected counts to one decimal place when needed.

\begin{enumerate}

  \item \textbf{Car Ownership Survey.} A random sample of 240 high school students
        is classified by grade (10th / 11th / 12th) and whether they own a car (Yes / No).
        The marginal totals are shown below.

        \vspace{0.1in}
        \renewcommand{\arraystretch}{1.6}
        \begin{tabular}{lcccc}
                     & \textbf{Yes} & \textbf{No} & \textbf{Row Total} \\
            \hline
            10th grade & \makebox[2cm][r]{\smash{\raisebox{-2pt}{\gradedblank[id=hw1a, type=short, points=1, answer=24, check=numeric, tolerance=0]{0.8cm}{0.18in}}}} & \makebox[2cm][r]{\smash{\raisebox{-2pt}{\gradedblank[id=hw1b, type=short, points=1, answer=36, check=numeric, tolerance=0]{0.8cm}{0.18in}}}} & 60 \\
            11th grade & \makebox[2cm][r]{\smash{\raisebox{-2pt}{\gradedblank[id=hw1c, type=short, points=1, answer=36, check=numeric, tolerance=0]{0.8cm}{0.18in}}}} & \makebox[2cm][r]{\smash{\raisebox{-2pt}{\gradedblank[id=hw1d, type=short, points=1, answer=54, check=numeric, tolerance=0]{0.8cm}{0.18in}}}} & 90 \\
            12th grade & \makebox[2cm][r]{\smash{\raisebox{-2pt}{\gradedblank[id=hw1e, type=short, points=1, answer=36, check=numeric, tolerance=0]{0.8cm}{0.18in}}}} & \makebox[2cm][r]{\smash{\raisebox{-2pt}{\gradedblank[id=hw1f, type=short, points=1, answer=54, check=numeric, tolerance=0]{0.8cm}{0.18in}}}} & 90 \\
            \hline
            Col Total  & 96 & 144 & 240 \\
        \end{tabular}

        \vspace{0.1in}
        \begin{teachernote}
            Sample work: 10th/Yes: $(60 \times 96)/240 = 24$. \quad
            11th/No: $(90 \times 144)/240 = 54$. \quad
            Notice 11th and 12th have equal row totals $\Rightarrow$ identical expected counts.
        \end{teachernote}

  \item \textbf{True or False.} \textcolor{keyred}{\textbf{FALSE}} \quad (TRUE)\gradedkey{\textit{T/F:}~\gradedblank[id=hw2a, type=choice, points=1, answer=F]{0.4in}{0.2in}}

        \vspace{0.1in}
        \ansline{False --- expected counts also depend on column totals. Equal row totals only guarantee that expected counts are proportional to column totals, not that they are all equal to each other.}\gradedkeyline{\gradedblank[id=hw2b, type=rubric, points=2, display={False: $E = (RT)(CT)/N$ also depends on the column totals; equal row totals make $E$ proportional to the column totals, not all equal.}]{\linewidth+6pt}{\gradedlinesheight{3}}}

  \item \textbf{Checking the Large-Counts Condition.} All expected counts exceed \ans{5}.\gradedkey{\gradedblank[id=hw3, type=short, points=1, answer=5, check=numeric, tolerance=0]{1cm}{0.2in}}

        \begin{enumerate}[label=\alph*.]
            \item Expected count for every cell:

                  \ansline{$E = \dfrac{(100)(75)}{300} = \ans{25}$ for every cell.}\gradedkey{$E = $~\gradedblank[id=hw3a, type=short, points=1, answer=25, check=numeric, tolerance=0]{0.9in}{0.2in}}

            \item \ansline{Yes --- $25 \ge 5$, so the large-counts condition is satisfied for all 12 cells.}\gradedkeyline{\gradedblank[id=hw3b, type=rubric, points=2, display={Yes: every expected count is $25 \ge 5$.}]{\linewidth+6pt}{\gradedlinesheight{1}}}
        \end{enumerate}

\end{enumerate}

\begin{reflectionbox}
\textbf{Reflection:} In your own words, explain why expected counts represent what we
would see ``if the two variables were independent.'' Use the idea of row proportions
and column totals in your answer.

\begin{teachernote}
    Model answer: If the variables are independent, the proportion of individuals in
    each row category is the same within every column. So the expected count in a cell
    is (row proportion)(column total) = (row total / table total)(column total). This
    equals the formula $E = (RT)(CT)/N$.
\end{teachernote}
\end{reflectionbox}

\end{document}
